\documentclass[11pt]{article}
\usepackage[T1]{fontenc}
\usepackage[margin=0.76in]{geometry}
\usepackage{amsmath,amssymb,amsthm}
\usepackage{booktabs,tabularx}
\usepackage{graphicx}
\usepackage[hidelinks]{hyperref}
\usepackage{microtype}

\newtheorem{theorem}{Theorem}
\newtheorem{proposition}[theorem]{Proposition}
\newtheorem{lemma}[theorem]{Lemma}
\newcommand{\F}{\mathbb F}
\newcommand{\VC}{\operatorname{VCdim}}
\newcommand{\sign}{\operatorname{sign}}
\newcommand{\qbinom}[2]{\genfrac{[}{]}{0pt}{}{#1}{#2}_2}

\title{The three-sparse Walsh threshold class:\\
an exact count and a one-bit VC-dimension gap}
\author{New partial result for Problem 8 of arXiv:1507.04319}
\date{13 August 2026}

\begin{document}
\maketitle

\paragraph{Status.}
\textbf{New partial result; the arbitrary-sparsity problem remains open.}
For one and two terms the VC dimension is exactly \(n+1\).  For three terms
the classifiers admit a complete algebraic description and exact count.  This
gives exact answers in dimensions \(n\le3\), and for every \(n\ge4\) it
determines the VC dimension up to one.

\section*{Source question}
Mixon and Peterson [1] define \(C_{n,k}\) to be the functions
\(g:\{\pm1\}^n\to\{\pm1\}\) of the form
\[
 g(x)=\sign\!\left(\sum_{i=1}^k a_i\prod_{j\in S_i}x_j\right)
\]
and ask in Problem 8: ``Determine the VC dimension of \(C_{n,k}\).''  They
prove
\(
\max\{n,k\}\le \VC(C_{n,k})\le2nk-O(k\log k).
\)

\begin{center}
\includegraphics[width=\textwidth]{figures/open_problem_crop.png}
\end{center}
\noindent\emph{Source evidence: arXiv:1507.04319, p.~7, Problem 8.}

We identify the cube with \(V=\F_2^n\), write
\(\chi_a(x)=(-1)^{a\cdot x}\), and encode a sign function as
\(g(x)=(-1)^{f(x)}\) with \(f:V\to\F_2\).  An affine Boolean form is a
function \(\ell(x)=b+a\cdot x\).

\section*{Proof intuition}
Two Walsh characters are rigid: after factoring one character, the sign of
\(a+b\chi\) is simply the sign of its larger coefficient, unless the
expression vanishes.  Three characters are the first genuinely new case.
After factoring one character they give a linear threshold function on a
two-dimensional Boolean square.  The only non-affine possibilities have one
exceptional vertex.  In algebraic normal form these are exactly
\(\ell+L_1L_2\).

The quadratic part of \(L_1L_2\) is the decomposable bivector
\(u\wedge v\), so the nonzero quadratic parts are indexed by two-planes in
\(\F_2^n\).  This gives the exact count and hence the upper bound.  For the
lower bound, start with \(0,e_1,\ldots,e_n\), then add
\(e_i+e_j\) for every edge of \(K_{2,n-2}\).  Affine terms fit arbitrary
labels on the affine basis.  Each right vertex of \(K_{2,n-2}\) then stores
its two remaining second differences as one vector in \(\F_2^2\), making all
edge labels independent.  The source problem for larger \(k\), and one bit
for \(k=3,n\ge4\), are not settled by this argument.

\begin{lemma}[two-term rigidity]\label{lem:two}
For every \(n\ge1\),
\[
 C_{n,2}=C_{n,1}=\{\pm\chi_a:a\in\F_2^n\},
 \qquad
 \VC(C_{n,1})=\VC(C_{n,2})=n+1.
\]
\end{lemma}

\begin{proof}
After combining repeated characters, a two-term expression with distinct
characters factors as
\[
 a\chi_r+b\chi_s=\chi_r(a+b\chi_{r+s}).
\]
If its sign is defined everywhere, then \(|a|\ne|b|\).  If
\(|a|>|b|\), its sign is \(\sign(a)\chi_r\); if \(|b|>|a|\), its sign is
\(\sign(b)\chi_s\).  Thus two terms give no new classifiers.

The \(2^{n+1}\) functions \(\pm\chi_a\) are precisely the phases of affine
Boolean forms.  They shatter \(\{0,e_1,\ldots,e_n\}\), since an affine form
can be prescribed arbitrarily on an affine basis.  The class-size bound gives
the reverse inequality.
\end{proof}

\begin{proposition}[classification and exact count]\label{prop:class}
The phases in \(C_{n,3}\) are exactly
\begin{equation}\label{eq:quadratic-class}
 (-1)^{\ell(x)+L_1(x)L_2(x)},
\end{equation}
where \(\ell,L_1,L_2:V\to\F_2\) are affine forms.  Consequently
\begin{equation}\label{eq:count}
 |C_{n,3}|=2^{n+1}\left(1+\qbinom{n}{2}\right)
 =2^{n+1}\left(1+\frac{(2^n-1)(2^n-2)}6\right),
\end{equation}
where the Gaussian coefficient is understood as zero for \(n<2\).
\end{proposition}

\begin{proof}
Consider three distinct characters and factor the first:
\[
 \sign(a\chi_r+b\chi_s+c\chi_t)
 =\chi_r\,\sign(a+b\chi_{r+s}+c\chi_{r+t}).
\]
The nonzero vectors \(r+s\) and \(r+t\) are distinct and hence linearly
independent over \(\F_2\).  The last sign is an affine real threshold
dichotomy on the four points of \(\{\pm1\}^2\).  Of the sixteen
dichotomies, parity and its complement are the only two that are not linearly
separable.  The remaining fourteen are constants, signed coordinates, or
dichotomies with exactly one exceptional vertex.  Constants and coordinates
give affine Boolean phases.  An exceptional-vertex dichotomy has the form
\((-1)^{L_1L_2}\), up to an affine phase.  Conversely,
\[
 (-1)^{pq}=\sign(1+(-1)^p+(-1)^q)
 \qquad(p,q\in\F_2),
\]
so every phase in \eqref{eq:quadratic-class} is the sign of at most three
characters.

Write the linear parts of \(L_1,L_2\) as \(u\cdot x,v\cdot x\).  The
quadratic part of \(L_1L_2\) is
\[
 \sum_{i<j}(u_iv_j+u_jv_i)x_ix_j,
\]
the bivector \(u\wedge v\).  It is zero when \(u,v\) are dependent.  If it
is nonzero, it uniquely determines the plane \(\langle u,v\rangle\): over
\(\F_2\), every ordered basis of that plane has determinant one and hence the
same nonzero wedge.  Thus the possible quadratic parts are zero and one for
each two-plane of \(V^*\).  For each quadratic part, the unique algebraic
normal form allows exactly \(2^{n+1}\) affine perturbations.  The number of
two-planes is \(\qbinom{n}{2}=(2^n-1)(2^n-2)/6\), proving
\eqref{eq:count}.
\end{proof}

\begin{theorem}[one-bit determination for three terms]\label{thm:main}
The three-sparse class satisfies
\[
 \VC(C_{1,3})=2,\qquad \VC(C_{2,3})=4,\qquad
 \VC(C_{3,3})=7.
\]
For every \(n\ge4\),
\begin{equation}\label{eq:one-bit}
 3n-3\le \VC(C_{n,3})\le3n-2.
\end{equation}
\end{theorem}

\begin{proof}
We first prove the lower bound for \(n\ge4\).  Let
\[
 S=\{0,e_1,\ldots,e_n\}
 \cup\{e_i+e_j:i\in\{1,2\},\ 3\le j\le n\}.
\]
Thus \(|S|=n+1+2(n-2)=3n-3\).  Prescribe arbitrary labels
\(y:S\to\F_2\).  For \(i\in\{1,2\}\), \(j\ge3\), set the requested second
difference
\[
 d_{ij}=y(0)+y(e_i)+y(e_j)+y(e_i+e_j).
\]
Choose \(u,v\in\F_2^n\) by
\[
 (u_1,v_1)=(1,0),\qquad (u_2,v_2)=(0,1),\qquad
 (u_j,v_j)=(d_{2j},d_{1j})\quad(j\ge3).
\]
For \(Q(x)=(u\cdot x)(v\cdot x)\), the coefficient of \(x_ix_j\) is
\(u_iv_j+u_jv_i=d_{ij}\) on every displayed edge.  Choose the unique affine
form \(\ell\) such that \(\ell+Q\) agrees with \(y\) on
\(0,e_1,\ldots,e_n\).  The four-point affine cancellation identity then
gives, for every added pair-sum point,
\[
 (\ell+Q)(0)+(\ell+Q)(e_i)+(\ell+Q)(e_j)
 +(\ell+Q)(e_i+e_j)=d_{ij},
\]
so \(\ell+Q=y\) throughout \(S\).  Proposition~\ref{prop:class} shows that
this is a member of \(C_{n,3}\).  Hence \(S\) is shattered.

For \(n=3\), take \(S=\{x\in\F_2^3:|x|\le2\}\).  Every alternating
quadratic form in three variables has rank at most two and is therefore
decomposable.  The same affine-interpolation argument realizes all labels on
these seven points.  Formula \eqref{eq:count} gives
\(|C_{3,3}|=128\), so the VC dimension is seven.  For \(n=2\), every Boolean
quadratic has decomposable quadratic part, so Proposition~\ref{prop:class}
gives all \(16\) functions; for \(n=1\), it gives all four.

Finally, for \(n\ge4\), \eqref{eq:count} implies
\[
 |C_{n,3}|
 =2^{n+1}\left(1+\frac{(2^n-1)(2^n-2)}6\right)
 <2^{3n-1}.
\]
A class shattering \(3n-1\) points would contain at least \(2^{3n-1}\)
distinct functions.  Therefore \(\VC(C_{n,3})\le3n-2\), completing
\eqref{eq:one-bit}.
\end{proof}

\small
\section*{Verification, limitations, and upgrade audit}
The proof is symbolic.  The bundled script independently checks
\(C_{n,2}=C_{n,1}\) for \(1\le n\le5\), the exact count for
\(1\le n\le5\), and every target labeling produced by the explicit
shattering synthesis for \(3\le n\le6\).  It also exhaustively computes
\(\VC(C_{n,3})=2,4,7,9\) for \(n=1,2,3,4\), consistent with the theorem.
These finite checks are sanity checks, not substitutes for the proof.

The remaining one-bit upgrade has a precise exterior-algebra formulation.
After quotienting restrictions by affine functions, shattering \(3n-2\)
points would require a \(2n-3\)-dimensional point-evaluation quotient of
\(\bigwedge^2\F_2^n\) in which every vector has a decomposable representative.
The construction above attains dimension \(2n-4\).  Direct quotient searches
for \(n=5,6\) found no final-dimension example, but no proof excluding one was
found; no exact claim is made for \(n\ge4\).

Exact-title, exact-question, Walsh--Hadamard VC-dimension, sparse polynomial
threshold, and threshold-density searches through 13 August 2026 found no
later source settling this exact problem or the formulas above.  The search
was bounded rather than a systematic bibliographic review.  Human review
should focus on the two-variable threshold classification, uniqueness in the
exact count, and the \(K_{2,n-2}\) second-difference synthesis.  The original
problem for arbitrary \(k\) remains open.

\subsection*{Reference}
\noindent [1] D. G. Mixon and J. Peterson,
\emph{Learning Boolean functions with concentrated spectra},
arXiv:1507.04319 (2015), especially Sections 2 and 6, Problem 8.

\end{document}
